% Reframed abstract and introduction for BrickForma.
% Original Xiao wording with the November 16 evaluation result statements updated.
% New citation keys to add to the .bib: legopuzzles (Tang et al., 2025, arXiv:2503.19990),
% pfp (Zhou et al., AAAI 2026, Paper Folding Puzzles), geogrambench (Luo et al., ICLR 2026).




\section{Introduction}
Spatial reasoning and geometric understanding are fundamental to human cognition and to artificial intelligence, supporting applications from robotics and navigation to computer-aided design.
As large language models (LLMs) have matured, a growing body of work asks whether they can reason about shape, position, and transformation. Recent benchmarks probe LLMs' ability through 
%multi-step LEGO assembly~\cite{brickgpt, bricknet},
paper folding and unfolding~\cite{zhou2026paperfolding} and geometry expressed as procedural drawing code~\cite{luo2026geogrambench}. Each reports pronounced deficiencies in current models. These studies largely assess spatial reasoning by asking a model to select or compute an answer. A complementary and less examined question is whether a language model can \emph{construct}: given only words, can it produce a three-dimensional object whose parts are correctly sized, correctly placed, and physically joined?

%xiao: use lego instead of Brick
Lego brick assembly makes this question concrete. A concise narrative request such as \emph{a simple chair} identifies an object category while leaving its dimensions, proportions, and part layout largely unspecified. Answering it requires three coupled forms of geometric spatial reasoning. The model must complete the missing geometry from prior knowledge of what the object looks like. It must discretize that geometry into pieces of fixed dimensions at integer positions on a grid. And it must respect the difference between spatial adjacency and stud connectivity, since pieces that touch side by side do not hold together unless they interlock vertically through studs. We refer to this setting as the \texttt{Text-to-Brick} task: the translation from a natural-language description, at any level of detail, to a set of (Lego) brick placements that reproduces the described shape as a single connected assembly. Because every output can be checked for coverage, bounds, and connectivity, the task provides a verifiable test of whether linguistic knowledge has become spatial structure.

Prior work has established that language models can be trained for this task. BrickGPT~\cite{brickgpt} learns autoregressive brick-placement sequences and incorporates validity checks with physics-aware rollback during inference. BrickNet~\cite{bricknet} adopts a graph-based representation supporting a broader range of parts and connection types. Two questions remain open alongside these advances. \textbf{(Q1)} How well do generators translate \emph{limited} linguistic information into object geometry? In our five-level benchmark (Table~\ref{tab:levels}), BrickNet obtains lower mean VQAScore, PE, and SigLIP2 than BrickForma at the two shortest description levels, L1 and L2, a gap that only level-wise evaluation makes visible. \textbf{(Q2)} Where should the geometric knowledge needed to answer concise requests come from? Large models hold broad knowledge of object proportions and part structure, but that knowledge is implicit and their outputs carry no structural guarantee. Whether it can be transferred to a compact model in a form that can be inspected and checked has received limited attention.

We present BrickForma, which addresses both questions by \emph{crystallizing} large-model geometric knowledge into a small model through executable programs rather than through the large model's direct outputs. With GPT-6.1 Sol assistance, we develop parametric programs that vary object proportions, part dimensions, openings, and support arrangements to synthesize voxel shapes in batches.
In this form, what a large model knows about a category becomes explicit and executable: each program fixes the structure of an object family and exposes its geometric degrees of freedom as parameters, so every generated shape is exact by construction. The source collection contains 80,000 models across 21 categories on a $48^3$ grid. Each model is paired with five levels of description that progressively introduce structural, geometric, proportional, and appearance information. The resulting approximately 400,000 caption records link concise requests and detailed specifications to shared geometry, so that a model trained on them learns to supply at the shortest level the geometric details that the longest level states outright.

Turning voxel shapes into brick targets is where spatial adjacency and stud connectivity must be separated. A connected voxel volume can be tiled by pieces that touch side by side without connecting through studs. Our conversion pipeline combines same-color tiling, connection analysis, and local bridging and retiling to construct geometrically checked target layouts, with changes to source voxels recorded in derived models. The integrated dataset release includes 79,301 models passing its dataset-level geometric checks. We fine-tune Qwen3-1.7B~\cite{qwen3} with low-rank adaptation~\cite{lora} on caption--placement pairs to predict brick dimensions and spatial positions. The knowledge transfer therefore passes through three inspectable stages, from programs to checked brick layouts to a 1.7B-parameter generator, and the small model never trains on an unverified large-model output.

Geometric reconstruction complements the small model's predictions with explicit structural enforcement. From the supported, in-bounds portions of predicted placements, the pipeline forms a voxel union and uses mixed-integer linear programming to seek an exact-cover layout forming a single stud-connected component. Independent coverage and connectivity checks verify the reconstructed result. This design assigns the two halves of the task to the components suited to them: the language model decides what shape a description denotes, and the solver decides how that shape holds together.
% TODO (optional, strengthens the framing): report the fine-tuned model's raw placements
% before reconstruction, and the un-tuned Qwen3-1.7B, to isolate what crystallization contributes.

We evaluate BrickForma against BrickGPT and BrickNet on 265 shared prompts spanning five description levels. BrickForma's raw outputs achieve the highest overall mean VQAScore, PE, and SigLIP2, with a mean invalid-placement ratio of 0.81\%. Reporting all five levels connects the multilevel data design to an explicit assessment of how generation changes with the amount of descriptive detail. At the shortest level, BrickForma leads PE and SigLIP2, while BrickGPT leads VQAScore.

In summary, the main contributions of this work are as follows:

\textbf{Crystallizing large-model geometric knowledge into a small model 
without human-verified data.}
We propose a data generation pipeline for transferring geometric and 
spatial knowledge from large language models to a compact model: the 
large model's knowledge is externalized as executable parametric programs, 
converted into geometrically validated training data, and distilled into 
a 1.7B-parameter model without relying on human verification. 
On the matched prompts, BrickForma's mean VQAScore is 16.5\% higher than
that of BrickGPT, the stronger baseline on this metric.



\textbf{A multi-level dataset that enables experimental validation of small language model alignment across 
various levels of descriptive detail.}
With our LLM-assisted parametric data construction pipeline, we generated 80,000 
source voxel models. Each voxel model is paired with five levels of textual description 
and geometrically validated brick targets. We fine-tuned Qwen3-1.7B on this dataset and carried out level-wise evaluation 
against BrickGPT and BrickNet to assess text--image alignment and invalid 
placements across descriptions ranging from concise prompts to detailed 
specifications. BrickForma achieves the highest mean PE and SigLIP2 across all five levels
and the highest mean VQAScore at L2--L5.